\chapter{Configuration and Grid Independence of the Topology Optimization Benchmarks}
\label{app:B_validation_configurations}

\providecommand{\meshplaceholder}[1]{%
	\fbox{\parbox[c][0.20\textheight][c]{0.88\linewidth}{\centering \textbf{Placeholder}\\[0.5em]#1}}%
}

This appendix collects the detailed setup information for the topology optimization benchmarks used in Chapter~\ref{ch:benchmark_results}. The pipe-bend benchmark, U-bend benchmark, and diffuser benchmark are included. For each case, the appendix records the geometry and physical parameters, the FEniCS simulation and optimization parameters, and the Ansys Fluent mesh-independence tables used for sharp-geometry CFD validation.

\section{Validation and Meshing Strategy}
\label{sec:app_B_mesh_strategy}
The FEniCS meshes used during topology optimization are treated as fixed optimization discretizations rather than as a separate mesh-convergence hierarchy. This distinction is important because refining the FEniCS mesh would also change the $DG_0$ design space, the filter stencil, the number of MMA design variables, and the diffuse Brinkman interface location. The pipe-bend optimization uses a rectangular-coordinate triangular mesh with near-wall clustering based on a smooth-wall $y^+\approx1$ estimate at the prescribed inlet Reynolds number. The U-bend optimization uses the corresponding wall-resolved triangular Gmsh mesh, with a near-wall size field based on the same $y^+\approx1$ estimate and a far-field target size of $h_{\max}=0.007L$ used to set the density-filter scale. The diffuser benchmark from Yoon~\cite{Yoon2016} uses a wall-resolved rectangular-coordinate triangular mesh over the square diffuser domain, refined around the full-height inlet, the center-third outlet, and the no-slip walls. In all cases, the mesh is chosen to provide sufficient wall-resolved SA and design resolution for the optimization problem; quantitative grid-convergence evidence is instead reported for the extracted sharp geometries in Ansys Fluent.

For each extracted turbulent topology optimization design, the Ansys validation mesh is refined through coarse, medium, and fine resolutions. Both SA and SST $k-\omega$ validation runs use the same extracted geometry, inlet condition, outlet condition, and fluid properties. The mesh-independence comparison is reported for both turbulence models because near-wall behavior and separated-flow predictions may respond differently to grid refinement.

All Fluent mesh-independence runs in this appendix use the same convergence and acceptance criteria. The static pressure drop must change by less than $5\times10^{-4}$ over the final $100$ iterations, while the viscous dissipation must change by less than $10^{-3}$ over the final $100$ iterations. The fine mesh is accepted only when the medium--fine relative difference is below $1\%$ for static pressure drop and below $2\%$ for viscous dissipation for both SA and SST $k-\omega$. The reported refinement differences are computed as $|\phi_{\mathrm{coarse}}-\phi_{\mathrm{medium}}|/|\phi_{\mathrm{coarse}}|$ and $|\phi_{\mathrm{medium}}-\phi_{\mathrm{fine}}|/|\phi_{\mathrm{fine}}|$ for each validation metric.

\section{Shared Turbulence-Topology Coupling Conventions}
\label{sec:app_B_shared_coupling_conventions}
All three turbulent topology optimization benchmarks use the same basic coupling idea between the density field, the SA working variable, and the reciprocal wall-distance equation. Topology-created solids act as walls for the SA working-variable penalty and the reciprocal wall-distance solve. The wall-density source is therefore \texttt{design}, the wall-distance mode is \texttt{reciprocal\_penalized}, and both the SA and wall-distance penalties use the power-law solid indicator defined in Chapter~\ref{ch:topology_optimization_methodology}. The same frozen-turbulence adjoint convention is also used: the forward loop updates $\tilde{\nu}$, $\nu_t$, and $G$, while the derivative freezes $\tilde{\nu}$ and $G$ during gradient assembly.

The pipe-bend and U-bend benchmarks then share a Bayat-style pressure-loss setup, while the diffuser follows the Yoon dissipation benchmark. The common mechanisms and the benchmark-specific parameter values are summarized in Table~\ref{tab:app_B_shared_coupling_conventions}.

The pipe-bend finite-difference sensitivity verification in Chapter~\ref{ch:benchmark_results} is not assigned a separate configuration table. It uses the same geometry, physical parameters, wall-distance convention, penalty settings, objective, filter, and solver setup as the pipe-bend topology optimization run. The only additional settings are the finite-difference sampling parameters listed in Table~\ref{tab:app_B_alexandersen_pipebend_numerical}.

\begin{table}[htbp]
	\centering
	\footnotesize
	\caption{Shared turbulence-topology coupling conventions and benchmark-specific parameter values.}
	\label{tab:app_B_shared_coupling_conventions}
	\renewcommand{\arraystretch}{1.2}
	\begin{tabularx}{\textwidth}{l X X}
		\hline
		\textbf{Parameter} & \textbf{Pipe bend and U-bend} & \textbf{Diffuser} \\
		\hline
		SA wall-distance mode & \texttt{reciprocal\_penalized} & \texttt{reciprocal\_penalized} \\
		SA wall-density source & \texttt{design} & \texttt{design} \\
		SA penalty interpolation & Power-law solid indicator & Power-law solid indicator \\
		Wall-distance penalty interpolation & Power-law solid indicator & Power-law solid indicator \\
		Wall-distance recovery & Reciprocal recovery from the relaxed Eikonal variable $G$ & Reciprocal recovery from the relaxed Eikonal variable $G$ \\
		Adjoint treatment & Frozen-turbulence adjoint; $\tilde{\nu}$ and $G$ frozen during gradient assembly & Frozen-turbulence adjoint; $\tilde{\nu}$ and $G$ frozen during gradient assembly \\
		Solid threshold for wall source & $\bar{\gamma}=0.50$ & No separate threshold parameter in the diffuser configuration \\
		SA working-variable penalty & Amplitude schedule ending at $10^5$; exponent $n=4$ & Amplitude $10^5$; exponent $n=3$ \\
		Wall-distance penalty & Amplitude schedule ending at $10^5$; exponent $n=4$ & Amplitude $10^5$; exponent $n=3$ \\
		Wall-distance parameters & $\sigma_w=0.01$, $G_0=20$ & $\sigma_w=0.1$, $G_0=20$ \\
		SA inlet treatment & Turbulence-intensity estimate: $I=7.5\%$, $\ell/H=0.05$, eddy-viscosity ratio capped at $50$ & Prescribed $\tilde{\nu}=1$ at inlet and outlet \\
		Optimization objective & Average inlet pressure with fixed outlet pressure $p=0$ & Dissipation objective $J_D$ with prescribed inlet and outlet velocity profiles \\
		\hline
	\end{tabularx}
\end{table}

\FloatBarrier

\section{Pipe Bend Benchmark}
\label{sec:app_B_config_alexandersen_pipebend}

The pipe-bend case follows the setup of Bayat et al.~\cite{Bayat2026}. The reference paper uses a standard $k$-$\varepsilon$ model with wall functions; the present thesis uses the wall-resolved SA implementation and evaluates the $V_f=0.25$ design.

\begin{table}[htbp]
	\centering
	\footnotesize
	\caption{Geometric, physical, and boundary-condition parameters for the pipe-bend benchmark.}
	\label{tab:app_B_alexandersen_pipebend_physical}
	\renewcommand{\arraystretch}{1.2}
	\begin{tabularx}{\textwidth}{l X}
		\hline
		\textbf{Parameter} & \textbf{Setting} \\
		\hline
		Reference case & Pipe-bend benchmark from Bayat et al.~\cite{Bayat2026} \\
		Reference figure for setup & Bayat et al. Fig.~10 \\
		Reference velocity figure & Bayat et al. Fig.~15 \\
		Reference turbulence model & Standard $k$-$\varepsilon$ with wall functions \\
		Turbulence model in this thesis & Spalart-Allmaras \\
		Design-domain dimensions & Square design region $L\times L$ with $L=1$ \\
		Full computational domain & $x/L\in[-0.2,1]$, $y/L\in[-0.2,1]$ \\
		Non-design inlet duct & Length $0.2L$, height $0.2L$, inlet over $y/L\in[0.7,0.9]$ \\
		Non-design outlet duct & Width $0.2L$, length $0.2L$, outlet over $x/L\in[0.7,0.9]$ on the lower boundary \\
		Volume fraction & $V_f=0.25$ \\
		Reynolds number & $Re=5{,}000$, based on $U_{\mathrm{in}}H_{\mathrm{in}}\rho/\mu$ with $H_{\mathrm{in}}=0.2L$ \\
		Density and viscosity & $\rho=1.0$, $\mu=4.0\times10^{-5}$ \\
		Inlet condition & Uniform inlet velocity $\mathbf{u}=(1,0)$; SA inlet estimated from the shared turbulence-intensity settings in Table~\ref{tab:app_B_shared_coupling_conventions} \\
		Outlet condition & Pressure outlet $p=0$ on the lower outlet, with the normal/tangential component constraint used by the FEniCS configuration \\
		Wall conditions & No-slip velocity and $\tilde{\nu}=0$ on external walls; topology-created solids enter through the design-density penalties \\
		Validation metrics & Average inlet pressure, pressure drop, $\Phi_D$, and relative changes between FEniCS-SA, Ansys-SA, and Ansys-SST $k-\omega$ \\
		\hline
	\end{tabularx}
\end{table}

\begin{table}[htbp]
	\centering
	\footnotesize
	\caption{Numerical FEniCS solver, topology-optimization, and sensitivity-verification parameters for the pipe-bend benchmark.}
	\label{tab:app_B_alexandersen_pipebend_numerical}
	\renewcommand{\arraystretch}{1.2}
	\begin{tabularx}{\textwidth}{l c X}
		\hline
		\textbf{Parameter} & \textbf{Symbol / option} & \textbf{Setting} \\
		\hline
		Flow finite elements & $(\mathbf{u},p)$ & Taylor--Hood $P_2/P_1$ continuous Galerkin elements \\
		SA working-variable space & $\tilde{\nu}$ & $P_1$ continuous Galerkin \\
		Design-density space & $\gamma$ & $DG_0$ cellwise constant density \\
		hh & -- & \texttt{mesh\_yplus1.xdmf}; $39{,}128$ vertices, $77{,}334$ triangles, $h_{\max}=0.007L$ \\
		Objective & $J_p$ & Average inlet pressure \\
		Initial design density & $\gamma_0$ & $V_f=0.25$, matched to the filtered volume \\
		Brinkman penalty & $\alpha(\bar{\gamma})$ & RAMP form with $\alpha_{\mathrm{fluid}}=0$, $\alpha_{\mathrm{solid}}=100$ \\
		RAMP continuation & $q$ & Code schedule $q=1/q_{\alpha}=[1/150,\ 1/75,\ 1/30,\ 1/15]$ \\
		Projection continuation & $\beta$ & $[4,\ 6,\ 9,\ 13]$ with threshold $\eta=0.50$ \\
		MMA move limits & -- & $[0.05,\ 0.04,\ 0.03,\ 0.02]$ \\
		Maximum optimization iterations & -- & Four continuation stages with $25$ MMA iterations each; stop if objective change is below $10^{-5}$ for five iterations \\
		Filter radius & $r_f$ & $4h_{\max}=0.028L$ \\
		Forward flow solver & -- & IPCS flow solves for frozen-SA Picard updates, followed by a warm-started SNES Newton line-search final flow solve \\
		Frozen NS-SA coupling & -- & Three Picard updates per design iteration with SA relaxation factor $0.20$ \\
		Linear solver & -- & MUMPS for SNES/state and adjoint solves; BiCGSTAB with ILU preconditioning for IPCS warm-start solves \\
		Adjoint treatment & -- & Frozen-turbulence adjoint in $(\mathbf{u}^*,p^*)$; $\tilde{\nu}$ and $G$ are frozen during gradient assembly \\
		Sensitivity verification samples & -- & Five active DG0 cells, seed $13$ \\
		Sensitivity verification step & $\Delta h$ & $10^{-6}$, central finite difference \\
		Sensitivity verification fields held fixed & -- & SA working variable $\tilde{\nu}^{0}$ and reciprocal wall-distance field $G^{0}$ \\
		\hline
	\end{tabularx}
\end{table}

The extracted turbulent pipe-bend geometry is meshed in Ansys Fluent with a patch-conforming triangular bulk mesh and wall-resolved inflation on all no-slip walls. The first-layer height is kept fixed from the smooth-wall $y^+\approx1$ estimate, while the mesh hierarchy refines the bulk face size, curvature angle, and inflation distribution.

\begin{table}[htbp]
	\centering
	\scriptsize
	\caption{Ansys Fluent mesh-independence table for the extracted turbulent pipe-bend benchmark geometry.}
	\label{tab:app_B_alexandersen_pipebend_mesh}
	\renewcommand{\arraystretch}{1.2}
	\setlength{\tabcolsep}{4pt}
	\begin{tabular}{@{}lccc@{}}
		\hline
		\textbf{Parameter} & \textbf{Coarse\_Ansys} & \textbf{Medium\_Ansys} & \textbf{Fine\_Ansys} \\
		\hline
		Mesh method & Triangles & Triangles & Triangles \\
		Bulk face element size [m] & $1.10 \times 10^{-2}$ & $5.00 \times 10^{-3}$ & $3.25 \times 10^{-3}$ \\
		Curvature angle & $15^\circ$ & $7^\circ$ & $4.5^\circ$ \\
		Near-wall treatment & Wall-resolved inflation & Wall-resolved inflation & Wall-resolved inflation \\
		First-layer height $\Delta y_1$ [m] & $6.00 \times 10^{-4}$ & $6.00 \times 10^{-4}$ & $6.00 \times 10^{-4}$ \\
		Inflation layers & $8$ & $14$ & $18$ \\
		Growth rate & $1.15$ & $1.10$ & $1.08$ \\
		Inflation thickness [m] & $8.24 \times 10^{-3}$ & $1.68 \times 10^{-2}$ & $2.25 \times 10^{-2}$ \\
		Nodes & $5{,}147$ & $21{,}343$ & $44{,}042$ \\
		Elements & $7{,}551$ & $32{,}780$ & $68{,}743$ \\
		Viscous dissipation $\Phi_D^{\mathrm{SA}}$ [W/m] & $0.0143$ & $0.0151$ & $0.0153$ \\
		Relative change $\Phi_D^{\mathrm{SA}}$, Coarse--Medium [\%] & -- & $5.38$ & -- \\
		Relative change $\Phi_D^{\mathrm{SA}}$, Medium--Fine [\%] & -- & -- & $1.32$ \\
		Viscous dissipation $\Phi_D^{\mathrm{SST}\ k\text{-}\omega}$ [W/m] & $0.0150$ & $0.0164$ & $0.0166$ \\
		Relative change $\Phi_D^{\mathrm{SST}\ k\text{-}\omega}$, Coarse--Medium [\%] & -- & $9.04$ & -- \\
		Relative change $\Phi_D^{\mathrm{SST}\ k\text{-}\omega}$, Medium--Fine [\%] & -- & -- & $1.22$ \\
		Static pressure drop $\Delta p^{\mathrm{SA}}$ [Pa] & $0.1423$ & $0.1449$ & $0.1457$ \\
		Relative change $\Delta p^{\mathrm{SA}}$, Coarse--Medium [\%] & -- & $1.84$ & -- \\
		Relative change $\Delta p^{\mathrm{SA}}$, Medium--Fine [\%] & -- & -- & $0.55$ \\
		Static pressure drop $\Delta p^{\mathrm{SST}\ k\text{-}\omega}$ [Pa] & $0.1481$ & $0.1516$ & $0.1522$ \\
		Relative change $\Delta p^{\mathrm{SST}\ k\text{-}\omega}$, Coarse--Medium [\%] & -- & $2.36$ & -- \\
		Relative change $\Delta p^{\mathrm{SST}\ k\text{-}\omega}$, Medium--Fine [\%] & -- & -- & $0.40$ \\
		Selected validation mesh & \multicolumn{3}{c}{Fine\_Ansys} \\
		\hline
	\end{tabular}
\end{table}

\FloatBarrier

\section{U-Bend Benchmark}
\label{sec:app_B_config_alexandersen_ubend}

The U-bend case follows Fig.~11 and Table~4 of Bayat et al.~\cite{Bayat2026}. It uses the same turbulent topology optimization framework as the pipe bend, but forces a $180^\circ$ return path around a fixed internal separator.

\begin{table}[htbp]
	\centering
	\footnotesize
	\caption{Geometric, physical, and boundary-condition parameters for the U-bend benchmark.}
	\label{tab:app_B_alexandersen_ubend_physical}
	\renewcommand{\arraystretch}{1.2}
	\begin{tabularx}{\textwidth}{l X}
		\hline
		\textbf{Parameter} & \textbf{Setting} \\
		\hline
		Reference case & U-bend benchmark from Bayat et al.~\cite{Bayat2026} \\
		Reference figure/table for setup & Bayat et al. Fig.~11 and Table~4 \\
		Reference turbulence model & Standard $k$-$\varepsilon$ with wall functions \\
		Turbulence model in this thesis & Spalart-Allmaras \\
		Design-domain dimensions & Square design region $L\times L$ with $L=1$ \\
		Full computational domain & $x/L\in[-0.2,1]$, $y/L\in[0,1]$ \\
		Port geometry & Two left-side ports of height $0.2L$; inlet over $y/L\in[0.55,0.75]$, outlet over $y/L\in[0.25,0.45]$ \\
		Fixed separator & Bar thickness $0.10L$, rounded-tip radius $0.05L$, extending from $x/L=-0.2$ to tip at $x/L=0.5$ \\
		Volume fraction & $V_f=0.27$ \\
		Reynolds number & $Re=5{,}000$, based on $U_{\mathrm{in}}H_{\mathrm{port}}\rho/\mu$ with $H_{\mathrm{port}}=0.2L$ \\
		Density and viscosity & $\rho=1.0$, $\mu=4.0\times10^{-5}$ \\
		Inlet condition & Uniform inlet velocity $\mathbf{u}=(1,0)$; SA inlet estimated from the shared turbulence-intensity settings in Table~\ref{tab:app_B_shared_coupling_conventions} \\
		Outlet condition & Pressure outlet $p=0$ on the lower left port, with the component constraint used by the FEniCS configuration \\
		Wall conditions & No-slip velocity and $\tilde{\nu}=0$ on external walls and the fixed separator; topology-created solids enter through the design-density penalties \\
		Validation metrics & Average inlet pressure, pressure drop, $\Phi_D$, and relative changes between FEniCS-SA, Ansys-SA, and Ansys-SST $k-\omega$ \\
		\hline
	\end{tabularx}
\end{table}

\begin{table}[htbp]
	\centering
	\footnotesize
	\caption{Numerical FEniCS solver and topology-optimization parameters for the U-bend benchmark.}
	\label{tab:app_B_alexandersen_ubend_numerical}
	\renewcommand{\arraystretch}{1.2}
	\begin{tabularx}{\textwidth}{l c X}
		\hline
		\textbf{Parameter} & \textbf{Symbol / option} & \textbf{Setting} \\
		\hline
		Flow finite elements & $(\mathbf{u},p)$ & Taylor--Hood $P_2/P_1$ continuous Galerkin elements \\
		SA working-variable space & $\tilde{\nu}$ & $P_1$ continuous Galerkin \\
		Design-density space & $\gamma$ & $DG_0$ cellwise constant density \\
		Reference FEniCS mesh & -- & \texttt{mesh\_yplus1.xdmf}; wall-resolved triangular Gmsh mesh with $45{,}280$ vertices, $85{,}780$ triangles, near-wall refinement based on $y^+\approx1$, and far-field $h_{\max}=0.007L$ \\
		Objective & $J_p$ & Average inlet pressure \\
		Initial design density & $\gamma_0$ & $V_f=0.27$, matched to the filtered volume \\
		Brinkman penalty & $\alpha(\bar{\gamma})$ & RAMP form with $\alpha_{\mathrm{fluid}}=0$, $\alpha_{\mathrm{solid}}=100$ \\
		RAMP continuation & $q$ & Code schedule $q=1/q_{\alpha}=[1/150,\ 1/75,\ 1/35,\ 1/12]$ \\
		Projection continuation & $\beta$ & $[8,\ 10,\ 18,\ 18]$ with threshold $\eta=0.50$ \\
		MMA move limits & -- & $[0.05,\ 0.04,\ 0.03,\ 0.02]$ \\
		Maximum optimization iterations & -- & Four continuation stages with $25$ MMA iterations each; stop if objective change is below $10^{-5}$ for five iterations \\
		Filter radius & $r_f$ & $4h_{\max}=0.028L$ \\
		Forward flow solver & -- & IPCS flow solves for frozen-SA Picard updates, followed by a warm-started SNES Newton line-search final flow solve \\
		Frozen NS-SA coupling & -- & Three Picard updates per design iteration with SA relaxation factor $0.20$ \\
		Linear solver & -- & MUMPS for SNES/state and adjoint solves; BiCGSTAB with ILU preconditioning for IPCS warm-start solves \\
		Adjoint treatment & -- & Frozen-turbulence adjoint in $(\mathbf{u}^*,p^*)$; $\tilde{\nu}$ and $G$ are frozen during gradient assembly \\
		\hline
	\end{tabularx}
\end{table}

\begin{table}[htbp]
	\centering
	\footnotesize
	\caption{Ansys Fluent mesh-independence table for the extracted turbulent U-bend benchmark geometry.}
	\label{tab:app_B_alexandersen_ubend_mesh}
	\renewcommand{\arraystretch}{1.2}
	\begin{tabularx}{\textwidth}{X c c c}
		\hline
		\textbf{Parameter} & \textbf{Coarse} & \textbf{Medium} & \textbf{Fine} \\
		\hline
		Bulk mesh type & [TBD] & [TBD] & [TBD] \\
		Near-wall treatment & [TBD] & [TBD] & [TBD] \\
		First-layer height $\Delta y_1$ [m] & [TBD] & [TBD] & [TBD] \\
		Inflation layers & [TBD] & [TBD] & [TBD] \\
		Growth rate & [TBD] & [TBD] & [TBD] \\
		Elements / cells & [TBD] & [TBD] & [TBD] \\
		Maximum $y^+$, SA & [TBD] & [TBD] & [TBD] \\
		Maximum $y^+$, SST $k-\omega$ & [TBD] & [TBD] & [TBD] \\
		$\Phi_D^{\mathrm{SA}}$ [W/m] & [TBD] & [TBD] & [TBD] \\
		$\Phi_D^{\mathrm{SST}\ k\text{-}\omega}$ [W/m] & [TBD] & [TBD] & [TBD] \\
		$\Delta p^{\mathrm{SA}}$ [Pa] & [TBD] & [TBD] & [TBD] \\
		$\Delta p^{\mathrm{SST}\ k\text{-}\omega}$ [Pa] & [TBD] & [TBD] & [TBD] \\
		Relative change, Medium--Fine [\%] & \multicolumn{3}{c}{[TBD]} \\
		Selected validation mesh & \multicolumn{3}{c}{[TBD]} \\
		\hline
	\end{tabularx}
\end{table}

\FloatBarrier

\section{Diffuser Benchmark}
\label{sec:app_B_config_yoon_diffuser}

The diffuser case follows Fig.~20--22 of Yoon~\cite{Yoon2016}. Unlike the pipe-bend and U-bend benchmarks, the diffuser prescribes parabolic velocity profiles at both the inlet and the outlet and uses the dissipation objective $J_D$. Its turbulence-density coupling still uses the shared \texttt{design}, \texttt{reciprocal\_penalized}, power-law, and frozen-adjoint conventions summarized in Table~\ref{tab:app_B_shared_coupling_conventions}.

\begin{table}[htbp]
	\centering
	\footnotesize
	\caption{Geometric, physical, and boundary-condition parameters for the diffuser benchmark.}
	\label{tab:app_B_yoon_diffuser_physical}
	\renewcommand{\arraystretch}{1.2}
	\begin{tabularx}{\textwidth}{l X}
		\hline
		\textbf{Parameter} & \textbf{Setting} \\
		\hline
		Reference case & Diffuser benchmark from Yoon~\cite{Yoon2016} \\
		Reference figures for setup and topology & Yoon Fig.~20--22 \\
		Reference turbulence model & Spalart-Allmaras \\
		Turbulence model in this thesis & Spalart-Allmaras with frozen-turbulence adjoint \\
		Design-domain dimensions & Square design region $L\times L$ with $L=1~\mathrm{m}$ \\
		Full computational domain & $x/L\in[0,1]$, $y/L\in[0,1]$ \\
		Inlet geometry & Exported Fluent inlet segment from $(1.16\times10^{-3},1.38\times10^{-2})$ to $(1.16\times10^{-3},0.979)$, with $H_{\mathrm{in}}=0.96558~\mathrm{m}$ \\
		Outlet geometry & Exported Fluent outlet segment from $(0.999,0.336)$ to $(0.999,0.670)$, with $H_{\mathrm{out}}=0.33333~\mathrm{m}$ \\
		Volume fraction & $V_f=0.30$ fluid, corresponding to Yoon's $70\%$ solid mass usage \\
		Reynolds number & $Re=3{,}000$, based on $U_{\max,\mathrm{in}}L\rho/\mu$ \\
		Density and viscosity & $\rho=1000~\mathrm{kg/m^3}$, $\mu=1.0~\mathrm{Pa\,s}$ \\
		Inlet condition & Shifted parabolic velocity profile with $U_{\max,\mathrm{in}}=3~\mathrm{m/s}$; $\tilde{\nu}=1$ for the FEniCS-SA solve \\
		Outlet condition & Shifted parabolic outflow profile with the maximum velocity set by incompressible flow-rate conservation; $\tilde{\nu}=1$ for the FEniCS-SA solve; pressure pinned at $(0,0)$ for gauge fixing \\
		Wall conditions & No-slip velocity and $\tilde{\nu}=0$ on the top and bottom walls and on the right boundary outside the outlet \\
		Brinkman solid resistance & $\alpha_{\mathrm{solid}}=10^9$, $\alpha_{\mathrm{fluid}}=0$ \\
		Optimization objective & Dissipation objective $J_D$, including physical viscous dissipation and Brinkman drag over the design domain \\
		Validation metrics & $J_D$, pressure drop, $\Phi_D$, and relative changes between FEniCS-SA, Ansys-SA, and Ansys-SST $k-\omega$ \\
		\hline
	\end{tabularx}
\end{table}

For the body-fitted Fluent evaluations, the parabolic inlet and outlet velocity profiles are shifted to the extracted inlet and outlet segments. The listed heights are used as the profile heights, while the lower vertices define the local profile origins. The inlet maximum velocity is kept at $3~\mathrm{m/s}$ so that the nominal Reynolds number remains $Re=3{,}000$; the outlet maximum velocity is adjusted by incompressible flow-rate conservation.

\begin{table}[htbp]
	\centering
	\footnotesize
	\caption{Numerical FEniCS solver and topology-optimization parameters for the diffuser benchmark.}
	\label{tab:app_B_yoon_diffuser_numerical}
	\renewcommand{\arraystretch}{1.2}
	\begin{tabularx}{\textwidth}{l c X}
		\hline
		\textbf{Parameter} & \textbf{Symbol / option} & \textbf{Setting} \\
		\hline
		Flow finite elements & $(\mathbf{u},p)$ & Taylor--Hood $P_2/P_1$ continuous Galerkin elements \\
		SA working-variable space & $\tilde{\nu}$ & $P_1$ continuous Galerkin \\
		Design-density space & $\gamma$ & $DG_0$ cellwise constant density \\
		Reference FEniCS mesh & -- & \texttt{mesh\_yoon\_yplus1.xdmf}; wall-resolved triangular mesh with $38{,}745$ vertices, $76{,}704$ triangles, far-field $h_{\max}=L/180=5.56\times10^{-3}L$, and first-cell width $1.82\times10^{-3}L$ \\
		Objective & $J_D$ & Dissipation objective over the design domain \\
		Initial design density & $\gamma_0$ & $V_f=0.30$, matched to the filtered volume \\
		Brinkman penalty & $\alpha(\bar{\gamma})$ & RAMP form with $\alpha_{\mathrm{fluid}}=0$, $\alpha_{\mathrm{solid}}=10^9$ \\
		RAMP continuation & $q$ & 12-stage schedule: 0.01, 0.02, 0.04, 0.08, 0.08, 0.12, 0.20, 0.35, 0.50, 0.75, 1.00, 1.00 \\
		Projection continuation & $\beta$ & 12-stage schedule: 1, 2, 4, 8, 8, 12, 16, 32, 64, 96, 128, 128, with threshold $\eta=0.50$ \\
		MMA move limits & -- & 12-stage schedule: 0.05, 0.04, 0.03, 0.02, 0.012, 0.010, 0.0075, 0.005, 0.0035, 0.0025, 0.0015, 0.0010 \\
		Maximum optimization iterations & -- & 12-stage schedule: 50, 60, 80, 100, 200, 180, 220, 260, 260, 220, 260, 360 MMA iterations; stop if objective change is below $10^{-6}$ for ten iterations \\
		Filter radius & $r_f$ & $4h_{\max}=2.22\times10^{-2}L$ \\
		SA working-variable penalty & -- & Power-law solid indicator with amplitude $10^5$ and exponent $n=3$ \\
		Wall-distance penalty & -- & Power-law solid indicator with amplitude $10^5$ and exponent $n=3$ \\
		Wall-distance equation & -- & Reciprocal penalized relaxed Eikonal equation with $\sigma_w=0.1$ and $G_0=20$ \\
		Forward flow solver & -- & IPCS flow solves for frozen-SA Picard updates, followed by a warm-started SNES Newton line-search final flow solve \\
		Frozen NS-SA coupling & -- & Three Picard updates per design iteration with SA relaxation factor $0.15$ \\
		Linear solver & -- & MUMPS for SNES/state and adjoint solves; BiCGSTAB with ILU preconditioning for IPCS warm-start solves \\
		Adjoint treatment & -- & Frozen-turbulence adjoint in $(\mathbf{u}^*,p^*)$; $\tilde{\nu}$ and $G$ are frozen during gradient assembly \\
		\hline
	\end{tabularx}
\end{table}

The extracted turbulent diffuser geometry is meshed with the same wall-resolved triangular strategy as the laminar diffuser comparison in Table~\ref{tab:app_C_yoon_diffuser_lam_mesh}, but with a distinct coarse--medium--fine hierarchy. The first-layer height is kept fixed at $\Delta y_1=1.50\times10^{-3}~\mathrm{m}$ to match the laminar-design diffuser meshes, while the bulk size, curvature angle, and inflation-layer distribution are refined independently for the turbulent topology.

\begin{table}[htbp]
	\centering
	\scriptsize
	\caption{Ansys Fluent mesh-independence table for the extracted turbulent diffuser benchmark geometry.}
	\label{tab:app_B_yoon_diffuser_mesh}
	\renewcommand{\arraystretch}{1.2}
	\setlength{\tabcolsep}{4pt}
	\begin{tabularx}{\textwidth}{X c c c}
		\hline
		\textbf{Parameter} & \textbf{Coarse\_Ansys} & \textbf{Medium\_Ansys} & \textbf{Fine\_Ansys} \\
		\hline
		Mesh method & Triangles & Triangles & Triangles \\
		Near-wall treatment & Wall-resolved inflation & Wall-resolved inflation & Wall-resolved inflation \\
		Bulk face element size [m] & $7.50 \times 10^{-3}$ & $5.00 \times 10^{-3}$ & $3.50 \times 10^{-3}$ \\
		Curvature angle & $9^\circ$ & $6^\circ$ & $4^\circ$ \\
		First-layer height $\Delta y_1$ [m] & $1.50 \times 10^{-3}$ & $1.50 \times 10^{-3}$ & $1.50 \times 10^{-3}$ \\
		Inflation layers & $10$ & $14$ & $18$ \\
		Growth rate & $1.11$ & $1.09$ & $1.07$ \\
		Inflation thickness [m] & $2.51 \times 10^{-2}$ & $3.90 \times 10^{-2}$ & $5.10 \times 10^{-2}$ \\
		Nodes & $8{,}104$ & [TBD] & [TBD] \\
		Elements & $12{,}330$ & [TBD] & [TBD] \\
		Viscous dissipation $\Phi_D^{\mathrm{SA}}$ [W/m] & [TBD] & [TBD] & [TBD] \\
		Relative change $\Phi_D^{\mathrm{SA}}$, Coarse--Medium [\%] & -- & [TBD] & -- \\
		Relative change $\Phi_D^{\mathrm{SA}}$, Medium--Fine [\%] & -- & -- & [TBD] \\
		Viscous dissipation $\Phi_D^{\mathrm{SST}\ k\text{-}\omega}$ [W/m] & [TBD] & [TBD] & [TBD] \\
		Relative change $\Phi_D^{\mathrm{SST}\ k\text{-}\omega}$, Coarse--Medium [\%] & -- & [TBD] & -- \\
		Relative change $\Phi_D^{\mathrm{SST}\ k\text{-}\omega}$, Medium--Fine [\%] & -- & -- & [TBD] \\
		Static pressure drop $\Delta p^{\mathrm{SA}}$ [Pa] & [TBD] & [TBD] & [TBD] \\
		Relative change $\Delta p^{\mathrm{SA}}$, Coarse--Medium [\%] & -- & [TBD] & -- \\
		Relative change $\Delta p^{\mathrm{SA}}$, Medium--Fine [\%] & -- & -- & [TBD] \\
		Static pressure drop $\Delta p^{\mathrm{SST}\ k\text{-}\omega}$ [Pa] & [TBD] & [TBD] & [TBD] \\
		Relative change $\Delta p^{\mathrm{SST}\ k\text{-}\omega}$, Coarse--Medium [\%] & -- & [TBD] & -- \\
		Relative change $\Delta p^{\mathrm{SST}\ k\text{-}\omega}$, Medium--Fine [\%] & -- & -- & [TBD] \\
		Selected validation mesh & \multicolumn{3}{c}{Fine\_Ansys, pending metric confirmation} \\
		\hline
	\end{tabularx}
\end{table}

\FloatBarrier
